\documentclass[11pt]{amsart}
\usepackage{amsmath,amssymb,amsthm,booktabs}
\usepackage[margin=1.15in]{geometry}
\usepackage{hyperref}
\newtheorem{theorem}{Theorem}
\newtheorem{proposition}[theorem]{Proposition}
\newtheorem{lemma}[theorem]{Lemma}
\newtheorem{corollary}[theorem]{Corollary}
\theoremstyle{definition}
\newtheorem{remark}[theorem]{Remark}
\newcommand{\Prob}{\mathbb{P}}
\newcommand{\E}{\mathbb{E}}
\newcommand{\Var}{\operatorname{Var}}
\newcommand{\as}{\alpha}

\title[The longest minimum-weight path in a complete graph]{An upper bound to second order for the longest
minimum-weight path in a complete graph}
\author{Yitzchak Shmalo}
\date{November 2026}

\begin{document}
\begin{abstract}
Give the edges of the complete graph on $n$ vertices independent exponential weights and let $M_n$ be the largest
number of edges of a minimum-weight path between two vertices. Addario-Berry, Broutin and Lugosi proved that
$M_n/\log n\to\as$ in probability, where $\as\approx3.5911$ solves $\as\log\as-\as=1$, with the bounds
$\as\log n-L\log\log n\le M_n\le\as\log n+\omega(n)$. We prove that
$\Prob(M_n\ge b_n+x)\le C\as^{-x}$ for all $x\ge0$, where $b_n=\as\log n-(\log\log n)/(2\log\as)$. Hence
$M_n-\as\log n\to-\infty$ in probability and $\E M_n\le b_n+O(1)$. We also find the asymptotic value of the
expected number of pairs of vertices joined by a minimum-weight path of at least $k$ edges when
$k=\as\log n+O(\log\log n)$, show that at $k=\lceil b_n\rceil$ the second moment of this number tends to infinity
although its mean stays bounded, and prove that
$\Prob(M_n=n-1)=n/(2\,(n-1)!)$ for every $n\ge2$. The lower side of the question of van der Hofstad on the
fluctuations of $M_n$ stays open.
\end{abstract}
\maketitle

\section{Introduction}
Let $K_n$ be the complete graph on the vertex set $[n]=\{1,\dots,n\}$ and give its edges independent weights with the
exponential distribution of mean one. Almost surely no two distinct paths with the same endpoints have the same weight, so for $i\ne j$ there is a
unique path of minimum weight from $i$ to $j$; we write $H_n(i,j)$ for its number of edges, the \emph{hopcount}
between $i$ and $j$. This is first-passage percolation on the complete graph, also called the stochastic mean-field
model of distance. The weights of the minimum-weight paths were described by Janson~\cite{Janson}: divided by
$\log n/n$, the weight between two fixed vertices, the largest weight from a fixed vertex and the largest weight over
all pairs converge in probability to $1$, $2$ and $3$; the fluctuations of the last quantity were found by Bhamidi and
van der Hofstad~\cite{BvdH}. Hopcounts for other edge-weight distributions on the complete graph are studied
in~\cite{EGHN1,EGHN2}. For exponential weights, the minimum-weight paths from a fixed vertex form a tree which, with
its vertices listed in the order of their distance from the root, is a random recursive tree
\cite{vdHHVM,SM}. Hence $H_n(1,2)$ is of order $\log n$, with Gaussian fluctuations of order
$\sqrt{\log n}$~\cite{vdHHVM,HVM}, and
\[
R_n=\max_{j}H_n(1,j)
\]
has the law of the height of a random recursive tree with $n$ vertices, and $R_n/\log n\to e$ in probability by
theorems of Devroye~\cite{Devroye} and Pittel~\cite{Pittel}.

The subject of this paper is the largest hopcount over all pairs,
\[
M_n=\max_{1\le i<j\le n}H_n(i,j).
\]
Let $\as\approx3.5911$ be the solution in $(e,\infty)$ of
\begin{equation}\label{eq:alpha}
\as\log\as-\as=1 .
\end{equation}
Janson~\cite{Janson} observed that the tail bound of Devroye for the height of a random recursive tree and a union
bound over the $n$ roots give $\Prob(M_n\ge\beta\log n)\to0$ for every $\beta>\as$, and asked whether $M_n/\log n$
converges to a constant in $[e,\as]$. Addario-Berry, Broutin and Lugosi~\cite{ABL} proved that it converges to $\as$. More precisely,
their Theorem~1 states that for some constant $L$ and every function $\omega(n)\to\infty$, with probability tending
to one,
\begin{equation}\label{eq:ABL}
\as\log n-L\log\log n\le M_n\le\as\log n+\omega(n),
\end{equation}
and that $\as\log n-L\log\log n\le\E M_n\le\as\log n+L$. In his Saint-Flour lectures, van der
Hofstad~\cite[Open Problem~3.1]{vdH} asks how the fluctuations of $R_n-e\log n$ and of $M_n-\as\log n$ behave. For
the root the answer is known: Addario-Berry and Ford~\cite{ABF} proved that the height of a random recursive tree is
$e\log n-\tfrac32\log\log n+O(1)$ in expectation, with exponential tails around the expectation, and finer results
on this height are in~\cite{PS1,PS2,ChenMa}. For $M_n$ the bounds~\eqref{eq:ABL} are the best available.

Our main result removes the gap between the upper bound in~\eqref{eq:ABL} and the scale $\log\log n$ of the lower
bound. Throughout, $\theta=\log\as$ and, for $n\ge3$,
\begin{equation}\label{eq:bn}
b_n=\as\log n-\frac{\log\log n}{2\theta},\qquad \frac1{2\theta}\approx0.3911 .
\end{equation}

\begin{theorem}\label{thm:upper}
There is a constant $C$ such that for all $n\ge3$ and all real $x\ge0$,
\[
\Prob\bigl(M_n\ge b_n+x\bigr)\le C\,\as^{-x}.
\]
Consequently $\E M_n\le b_n+O(1)$, and $\Prob(M_n\ge b_n+w_n)\to0$ whenever $w_n\to\infty$. In particular
$M_n-\as\log n\to-\infty$ in probability.
\end{theorem}

Together with the lower bound in~\eqref{eq:ABL}, Theorem~\ref{thm:upper} shows that, with probability tending to
one,
\[
-L\log\log n\ \le\ M_n-\as\log n\ \le\ -\frac{\log\log n}{2\theta}+w_n
\]
for every sequence $w_n\to\infty$.

The proof is a first-moment computation in which the factor $(\log n)^{-1/2}$ of the local limit theorem is kept. For
$k\ge1$ let
\[
N_{n,k}=\#\bigl\{\{i,j\}:\ i\ne j,\ H_n(i,j)\ge k\bigr\}
\]
be the number of pairs at hopcount at least $k$, so that $\{M_n\ge k\}=\{N_{n,k}\ge1\}$. The mean of $N_{n,k}$ is
given exactly by the expected profile of the random recursive tree (Lemma~\ref{lem:first}), and its asymptotic
behaviour near $b_n$ is as follows.

\begin{theorem}\label{thm:first}
Let $k_n$ be integers with $k_n-\as\log n=O(\log\log n)$. Then
\[
\E N_{n,k_n}=\bigl(C_*+o(1)\bigr)\,\as^{-(k_n-b_n)},\qquad
C_*=\frac{1}{2\,\Gamma(1+\as)\,(1-1/\as)\sqrt{2\pi\as}}\approx0.01104 .
\]
\end{theorem}

Thus $b_n$ is the threshold of the first moment: $\E N_{n,k}$ is of order one exactly when $k=b_n+O(1)$. Whether
$M_n-b_n$ is tight is not known. For the root, the same computation puts the first-moment threshold of the height at
$e\log n-\frac12\log\log n$, whereas the expected height is $e\log n-\frac32\log\log n+O(1)$ by~\cite{ABF}, as for the maximum
of a branching random walk. The level $\as\log n$, however, lies in the large deviation range of the height of a
single tree, where the probability of reaching a level and the expected number of vertices at it differ at most by a
power of $\log n$~\cite{ChenMa}, so the comparison with the root does not predict the centring of $M_n$. What a
second moment argument meets is made precise by the next statement: at the first-moment threshold the pairs at large
hopcount come in clusters with a common endpoint, and the second moment of their number is unbounded.

\begin{proposition}\label{prop:second}
There is a constant $c>0$ such that, with $k_n=\lceil b_n\rceil$ and for all large $n$,
\[
\E\bigl[N_{n,k_n}^2\bigr]\ \ge\ c\,n^{\,5-\as\log4}\,(\log n)^{\frac{\log 4}{2\theta}-\frac32},
\qquad 5-\as\log4\approx0.0216>0 .
\]
The pairs $\{i,j\}$, $\{i,j'\}$ with a common endpoint alone contribute this much. In particular
$(\E N_{n,k_n})^2/\E[N_{n,k_n}^2]\to0$.
\end{proposition}

So the second moment method applied to $N_{n,k}$ itself gives nothing at $k=b_n+O(1)$; the lower bound
in~\eqref{eq:ABL} is obtained in~\cite{ABL} by applying it to a restricted family of paths, which leads to the
constant $L$. Proposition~\ref{prop:second} does not decide whether $\Prob(M_n\ge b_n-A)$ is bounded away from zero for
large $A$.

At the other end of the range the law of $M_n$ can be written down.

\begin{proposition}\label{prop:top}
For every $n\ge2$,
\[
\Prob(M_n=n-1)=\frac{n}{2\,(n-1)!}.
\]
\end{proposition}

For $n=3,4,5,6$ this gives $3/4$, $1/3$, $5/48$ and $1/40$.

The identity for $\E N_{n,k}$ and the tilting that leads to Theorems~\ref{thm:upper} and~\ref{thm:first} are in
Section~\ref{sec:profile}, the two theorems are proved in Sections~\ref{sec:upper} and~\ref{sec:local},
Proposition~\ref{prop:second} in Section~\ref{sec:second} and Proposition~\ref{prop:top} in Section~\ref{sec:top},
which also reports simulations of $M_n$.

\section{The expected number of pairs at large hopcount}\label{sec:profile}

Fix a root $r\in[n]$. The union of the minimum-weight paths from $r$ to the other vertices is a spanning tree, the
shortest-path tree of $r$. Explore it by increasing distance from $r$. When $m$ vertices have been found, the
memoryless property of the exponential distribution makes the residual weights of the $m(n-m)$ edges between found
and unfound vertices, that is, the weight of each such edge less the time elapsed since its found endpoint was
reached, independent exponential variables of mean one, so the next vertex is attached by an edge chosen
uniformly among them, and its parent is uniform among the $m$ vertices found so far, independently of the past. The
tree, with its vertices numbered in the order in which they are found, is therefore a random recursive tree: vertex
$m+1$ is attached to a uniform vertex among $1,\dots,m$~\cite{vdHHVM,SM}. The hopcount $H_n(r,v)$ is the depth of $v$
in this tree. The trees of different roots are dependent, and nothing below uses more than the law of one of them.

Let $Z_{n,d}$ be the number of vertices at depth $d$ in a random recursive tree with $n$ vertices, and
\[
F_n(z)=\E\sum_{d\ge0}Z_{n,d}\,z^d,\qquad f_{n,d}=[z^d]F_n(z)=\E Z_{n,d}.
\]
A new vertex attached to a uniform vertex among $n$ has depth one more than its parent, so
$F_{n+1}(z)=(1+z/n)F_n(z)$ and
\begin{equation}\label{eq:F}
F_n(z)=\prod_{j=1}^{n-1}\Bigl(1+\frac zj\Bigr)=\frac{\Gamma(n+z)}{\Gamma(n)\,\Gamma(1+z)} .
\end{equation}
Thus $(n-1)!\,f_{n,d}$ is the unsigned Stirling number of the first kind $\genfrac{[}{]}{0pt}{}{n}{d+1}$; this description of the expected profile is
classical, see~\cite{SM}, \cite[Chapter~6]{Drmota} and~\cite{FHN}.

\begin{lemma}\label{lem:first}
For all $n\ge2$ and $k\ge1$,
\[
\E N_{n,k}=\frac n2\sum_{d\ge k}f_{n,d}.
\]
\end{lemma}

\begin{proof}
Each pair $\{i,j\}$ is counted once from the root $i$ and once from the root $j$, so
$2N_{n,k}=\sum_{r}\#\{v\ne r:H_n(r,v)\ge k\}$. For each $r$ the inner number is the number of vertices at depth at
least $k$ in the shortest-path tree of $r$, whose mean is $\sum_{d\ge k}f_{n,d}$.
\end{proof}

The sum in Lemma~\ref{lem:first} is over the whole profile. It should not be replaced by $n-1$ times the tail of the
hopcount between two fixed vertices computed from $F_n(z)/n$: that function is the generating function of the depth
of a uniform vertex with the root included, and the depth of a uniform vertex other than the root has generating
function $(F_n(z)-1)/(n-1)$.

\subsection*{Tilting}
Let $a>0$. Let $B_1,\dots,B_{n-1}$ be independent Bernoulli variables with
\[
\Prob(B_j=1)=q_j=\frac{a}{j+a},\qquad S_n=\sum_{j=1}^{n-1}B_j .
\]
Since $(1+au/j)/(1+a/j)=1-q_j+q_ju$, we have $\E u^{S_n}=F_n(au)/F_n(a)$ and therefore
\begin{equation}\label{eq:tilt}
f_{n,d}=F_n(a)\,a^{-d}\,\Prob(S_n=d)\qquad(d\ge0).
\end{equation}
For $a=1$ this is the representation of the depth of a uniform vertex of the random recursive tree as a sum of
independent indicators~\cite{Devroye88}.
The mean and variance of $S_n$ are
\[
\mu_n=\sum_{j<n}\frac{a}{j+a},\qquad V_n=\sum_{j<n}\frac{aj}{(j+a)^2} .
\]
From $0\le a/j-a/(j+a)\le a^2/j^2$ and $|aj/(j+a)^2-a/j|\le 2a^2/j^2+a^3/j^3$,
\begin{equation}\label{eq:muV}
\mu_n=a\log n+O(1),\qquad V_n=a\log n+O(1),
\end{equation}
with constants depending on $a$ only. As $V_n>0$ for every $n\ge2$, there is $c_a>0$ with $V_n\ge c_a\log n$ for all
$n\ge2$.

\begin{lemma}\label{lem:atom}
Let $X$ be a sum of finitely many independent Bernoulli variables, with variance $V>0$. Then for all real $t$,
\begin{equation}\label{eq:cf}
\bigl|\E e^{itX}\bigr|\le\exp\bigl(-2V\sin^2(t/2)\bigr),
\end{equation}
and $\sup_{d}\Prob(X=d)\le\sqrt{\pi/(8V)}$. In particular, for all $n\ge2$,
\begin{equation}\label{eq:atom}
\sup_{d}\Prob(S_n=d)\le\sqrt{\frac{\pi}{8V_n}}\le\frac{C_a}{\sqrt{\log n}} .
\end{equation}
\end{lemma}

\begin{proof}
For one Bernoulli variable with parameter $q$, $|1-q+qe^{it}|^2=1-4q(1-q)\sin^2(t/2)$, and
$\sqrt{1-y}\le e^{-y/2}$ for $0\le y\le1$; multiplying over the summands gives~\eqref{eq:cf}. By Fourier inversion
for an integer-valued variable,
\[
\Prob(X=d)=\frac1{2\pi}\int_{-\pi}^{\pi}\E e^{itX}\,e^{-itd}\,dt
\le\frac1{2\pi}\int_{-\pi}^{\pi}\bigl|\E e^{itX}\bigr|\,dt
\le\frac1{2\pi}\int_{-\infty}^{\infty}e^{-2Vt^2/\pi^2}\,dt=\sqrt{\frac{\pi}{8V}},
\]
where we used $\sin(|t|/2)\ge|t|/\pi$ for $|t|\le\pi$.
\end{proof}

\section{Proof of Theorem~\ref{thm:upper}}\label{sec:upper}

From now on $a=\as$ in the tilting of Section~\ref{sec:profile}, and $\theta=\log\as$.

\begin{lemma}\label{lem:finite}
There is a constant $C$ such that for all $n\ge2$ and all integers $k\ge1$,
\[
\Prob(M_n\ge k)\le\E N_{n,k}\le C\,\frac{n^{1+\as}\,\as^{-k}}{\sqrt{\log n}} .
\]
\end{lemma}

\begin{proof}
The first inequality is Markov's. By Lemma~\ref{lem:first}, \eqref{eq:tilt} and~\eqref{eq:atom},
\[
\E N_{n,k}=\frac n2F_n(\as)\sum_{d\ge k}\as^{-d}\,\Prob(S_n=d)
\le\frac n2F_n(\as)\,\frac{\as^{-k}}{1-1/\as}\,\frac{C_{\as}}{\sqrt{\log n}} ,
\]
and $\log F_n(\as)\le\as\sum_{j<n}1/j\le\as(1+\log n)$.
\end{proof}

\begin{proof}[Proof of Theorem~\ref{thm:upper}]
Let $n\ge3$, write $\ell=\log n$, and let $x\ge0$. Since $M_n$ is an integer, $M_n\ge b_n+x$ exactly when
$M_n\ge k$ with $k=\lceil b_n+x\rceil\ge1$. By~\eqref{eq:alpha}, $1+\as-\as\theta=0$, so
\[
\frac{n^{1+\as}\,\as^{-b_n}}{\sqrt{\ell}}
=\exp\Bigl((1+\as-\as\theta)\,\ell+\tfrac12\log\ell\Bigr)\,\ell^{-1/2}=1 ,
\]
and Lemma~\ref{lem:finite} gives $\Prob(M_n\ge b_n+x)\le C\as^{-(k-b_n)}\le C\as^{-x}$. For the expectation, with
$m=\lceil b_n\rceil$,
\[
\E M_n\le m+\sum_{j\ge1}\Prob(M_n\ge m+j)\le b_n+1+C\sum_{j\ge1}\as^{-j} .
\]
The remaining statements follow by taking $x=w_n$ and $x=(\log\log n)/(2\theta)-A$.
\end{proof}

\begin{remark}
The bound of~\cite{ABL} for the upper side of~\eqref{eq:ABL} uses the inequality
$\Prob(\text{height}\ge x\log m)\le e^{x-1}m^{x-x\log x}$, $x>1$, of Devroye~\cite{Devroye}, which is the Chernoff bound for
the sum in Lemma~\ref{lem:first}. Lemma~\ref{lem:finite} differs from it by the factor $(\log n)^{-1/2}$ of
Lemma~\ref{lem:atom}, and this factor is the source of the term $-(\log\log n)/(2\theta)$. For the height $R_n$ of
one tree and a fixed $\beta>e$, Chen and Ma~\cite[Proposition~3.1]{ChenMa} prove
$\Prob(R_n\ge\beta\log n)\le Cn^{-\beta(\log\beta-1)}(\log n)^{-1/2}$, together with a lower bound of the same power
of $n$. With $\beta=\as$ and a union bound over the roots their inequality gives
$\Prob(M_n\ge\as\log n)=O((\log n)^{-1/2})$, which is the case $k=\lceil\as\log n\rceil$ of
Lemma~\ref{lem:finite}. The second-order term requires the bound at the levels $k=\as\log n-O(\log\log n)$, with a
constant that does not depend on the level, and this is what Lemma~\ref{lem:finite} provides.
\end{remark}

\section{Proof of Theorem~\ref{thm:first}}\label{sec:local}

We keep $a=\as$. The first step is a local limit theorem for sums of independent Bernoulli variables. It is
standard; we include the short argument because it is used for $S_n$ here and for two other families of sums in
Section~\ref{sec:second}.

\begin{lemma}\label{lem:local}
For $m\ge1$ let $X_m$ be a sum of finitely many independent Bernoulli variables, with mean $\mu^{(m)}$ and variance
$V^{(m)}$, and let $d_m$ be an integer. If $V^{(m)}\to\infty$ and $d_m-\mu^{(m)}=o\bigl(\sqrt{V^{(m)}}\bigr)$, then
$\sqrt{V^{(m)}}\,\Prob(X_m=d_m)\to1/\sqrt{2\pi}$.
\end{lemma}

\begin{proof}
Let $B$ be a Bernoulli variable with parameter $q$ and $v=q(1-q)$. The centred variable $B-q$ has variance $v$ and
absolute third moment at most $v$, so for $|t|\le1$,
\[
\E e^{it(B-q)}=1-\tfrac12vt^2+O(v|t|^3),\qquad
\log\E e^{it(B-q)}=-\tfrac12vt^2+O\bigl(v|t|^3\bigr),
\]
with absolute constants. Summing over the summands of $X_m$, we get
$\log\E\exp(it(X_m-\mu^{(m)}))=-\frac12V^{(m)}t^2+O(V^{(m)}|t|^3)$ for $|t|\le1$, and hence
$\E\exp(iu(X_m-\mu^{(m)})/\sqrt{V^{(m)}})\to e^{-u^2/2}$ for every real $u$. By Fourier inversion and the
substitution $t=u/\sqrt{V^{(m)}}$,
\[
\sqrt{V^{(m)}}\,\Prob(X_m=d_m)=\frac1{2\pi}\int_{-\pi\sqrt{V^{(m)}}}^{\pi\sqrt{V^{(m)}}}
\E e^{iu(X_m-\mu^{(m)})/\sqrt{V^{(m)}}}\;e^{-iu(d_m-\mu^{(m)})/\sqrt{V^{(m)}}}\,du .
\]
The integrand, set equal to zero outside the interval of integration, converges pointwise to $e^{-u^2/2}$ and, by~\eqref{eq:cf} and $\sin(|t|/2)\ge|t|/\pi$, its modulus is
at most $e^{-2u^2/\pi^2}$. Dominated convergence gives the limit $(2\pi)^{-1}\int e^{-u^2/2}du=1/\sqrt{2\pi}$.
\end{proof}

\begin{proof}[Proof of Theorem~\ref{thm:first}]
Write $\ell=\log n$. By~\eqref{eq:muV}, $\mu_n=\as\ell+O(1)$ and $V_n=\as\ell+O(1)$, so for each fixed integer
$r\ge0$ the point $k_n+r$ satisfies the hypothesis of Lemma~\ref{lem:local} and
\[
\sqrt{\ell}\;\Prob(S_n=k_n+r)\to\frac1{\sqrt{2\pi\as}} .
\]
The left side is bounded uniformly in $r$ and $n$ by~\eqref{eq:atom}, and $\sum_r\as^{-r}<\infty$, so by dominated
convergence
\[
\sum_{r\ge0}\as^{-r}\,\Prob(S_n=k_n+r)=\frac{1+o(1)}{(1-1/\as)\sqrt{2\pi\as\,\ell}} .
\]
By~\eqref{eq:F} and Stirling's formula, $F_n(\as)=(1+o(1))\,n^{\as}/\Gamma(1+\as)$. Lemma~\ref{lem:first}
and~\eqref{eq:tilt} now give
\[
\E N_{n,k_n}=\frac n2F_n(\as)\,\as^{-k_n}\sum_{r\ge0}\as^{-r}\,\Prob(S_n=k_n+r)
=\frac{(1+o(1))\;n^{1+\as}\,\as^{-k_n}}{2\,\Gamma(1+\as)(1-1/\as)\sqrt{2\pi\as\,\ell}} ,
\]
and $n^{1+\as}\as^{-b_n}/\sqrt{\ell}=1$, as in the proof of Theorem~\ref{thm:upper}.
\end{proof}

\begin{remark}
Theorem~\ref{thm:first} is a statement about sums of Stirling numbers of the first kind in the range where the
second index is close to $\as$ times the logarithm of the first. Asymptotic expansions of these numbers go back to
Moser and Wyman~\cite{MW}; the expansion of Hwang~\cite{Hwang} is uniform for a second index up to any fixed multiple
of the logarithm of the first, and Theorem~\ref{thm:first} could also be derived from it.
\end{remark}

\begin{remark}
The convergence in Theorem~\ref{thm:first} is slow. For $a=\as$ one has
$\mu_n=\as\log n-\as\psi(1+\as)+o(1)$, where $\psi$ is the digamma function and $\as\psi(1+\as)\approx5.07$, so at
$k=b_n+O(1)$ the Gaussian factor $\exp(-(k-\mu_n)^2/(2V_n))$ of the local limit theorem differs from one by a term
of order $(\log\log n)^2/\log n$, and for moderate $n$ the constant $\as\psi(1+\as)$ is the larger part of
$k-\mu_n$. At $k=\lceil b_n\rceil$ the ratio of $\E N_{n,k}$, computed exactly
from~\eqref{eq:F}, to $C_*\as^{-(k-b_n)}$ is $0.388$, $0.458$, $0.614$ and $0.718$ for $n=10^2,10^3,10^4,10^5$; the
Gaussian factor at the same points is $0.320$, $0.436$, $0.606$ and $0.714$.
\end{remark}

\section{The second moment at the first-moment threshold}\label{sec:second}

Let $n\ge4$ and $k\ge1$, and let $Y$ be the number of vertices at depth at least $k$ in a random recursive tree with
$n$ vertices. Expanding the square of $N_{n,k}$ as a sum over two pairs, which coincide, share one vertex or are
disjoint, and using the exchangeability of the vertices,
\begin{equation}\label{eq:square}
\E\bigl[N_{n,k}^2\bigr]=\binom n2p+6\binom n3q_s+6\binom n4q_d ,
\end{equation}
where $p=\Prob(H_n(1,2)\ge k)$, $q_s=\Prob(H_n(1,2)\ge k,\ H_n(1,3)\ge k)$ and
$q_d=\Prob(H_n(1,2)\ge k,\ H_n(3,4)\ge k)$. The first two probabilities are determined by the shortest-path tree of
vertex $1$: the $n-1$ other vertices are exchangeable, so
\[
p=\frac{\E Y}{n-1},\qquad q_s=\frac{\E[Y(Y-1)]}{(n-1)(n-2)},\qquad 6\binom n3q_s=n\,\E[Y(Y-1)] .
\]
The probability $q_d$ involves two trees and is not determined by the law of one of them. We bound the middle term
of~\eqref{eq:square} from below.

For a random recursive tree $T_j$ with $j$ vertices let
\[
R_j(z,w)=\E\sum_{u\ne v}z^{\operatorname{depth}(u)}\,w^{\operatorname{depth}(v)},
\]
the sum being over ordered pairs of distinct vertices of $T_j$.

\begin{lemma}\label{lem:pairs}
For all $n\ge2$,
\begin{equation}\label{eq:R}
R_n(z,w)=\sum_{j=1}^{n-1}\frac{z+w}{j}\,F_j(zw)\prod_{l=j+1}^{n-1}\Bigl(1+\frac{z+w}{l}\Bigr).
\end{equation}
\end{lemma}

\begin{proof}
Write $P_j(z)=\sum_{u\in T_j}z^{\operatorname{depth}(u)}$. If vertex $j+1$ is attached to a vertex of depth $D$, the
new ordered pairs contribute $z^{D+1}P_j(w)+w^{D+1}P_j(z)$. Given $T_j$, the parent is uniform, so the conditional
mean of this is $j^{-1}(z+w)P_j(z)P_j(w)$, and $P_j(z)P_j(w)$ is the sum over ordered pairs of distinct vertices
plus $P_j(zw)$. Hence
\[
R_{j+1}=\Bigl(1+\frac{z+w}{j}\Bigr)R_j+\frac{z+w}{j}\,F_j(zw),\qquad R_1=0,
\]
and~\eqref{eq:R} follows by induction.
\end{proof}

In the $j$th term of~\eqref{eq:R} the factor $F_j(zw)$ accounts for the common part of the two paths from the root,
which ends at the parent of vertex $j+1$, and the remaining factors account for the two separate parts. It is this
term, for $j$ a suitable power of $n$, that is large.

\begin{proof}[Proof of Proposition~\ref{prop:second}]
Let $k=k_n=\lceil b_n\rceil$ and $T=\log n$. Write $Q_{j,n}(t)=\prod_{l=j+1}^{n-1}(1+t/l)$ and
$q_{j,h}=[t^h]Q_{j,n}(t)$. In the $j$th term of~\eqref{eq:R} the monomial $z^kw^k$ arises from the term $(zw)^d$ of
$F_j(zw)$ and the term $z^sw^s$ of $(z+w)\,Q_{j,n}(z+w)$, where $s=k-d\ge1$; the latter has coefficient
$\binom{2s}{s}q_{j,2s-1}$. Therefore
\begin{equation}\label{eq:coef}
[z^kw^k]R_n=\sum_{j=1}^{n-1}\frac1j\sum_{d<k}f_{j,d}\;q_{j,2(k-d)-1}\binom{2(k-d)}{k-d}.
\end{equation}
All terms are nonnegative. Let $D_j$ be the sum of independent Bernoulli variables with parameters $4/(l+4)$,
$1\le l<j$, and $L_j$ the sum of independent Bernoulli variables with parameters $4/(l+4)$, $j<l<n$. As
in~\eqref{eq:tilt},
\[
f_{j,d}=F_j(4)\,4^{-d}\,\Prob(D_j=d),\qquad q_{j,h}=Q_{j,n}(4)\,4^{-h}\,\Prob(L_j=h),
\]
and $F_j(4)\,Q_{j,n}(4)=F_n(4)/(1+4/j)$. With $s=k-d$, \eqref{eq:coef} becomes
\begin{equation}\label{eq:coef2}
[z^kw^k]R_n=F_n(4)\,4^{1-k}\sum_{j=1}^{n-1}\frac1{j+4}\sum_{d<k}\Prob(D_j=d)\,\Prob(L_j=2s-1)\,
\binom{2s}{s}4^{-s}.
\end{equation}

Put $\rho=(\as-2)/2\in(0,1)$ and $N=\lfloor n^{\rho}\rfloor$. We keep in~\eqref{eq:coef2} only the indices
$N\le j\le2N$ and, for each of them, the single value $d=d_j=\lfloor\E D_j\rfloor$. As in~\eqref{eq:muV}, uniformly
in these $j$,
\[
\E D_j=4\log j+O(1)=4\rho T+O(1),\qquad \E L_j=4\log(n/j)+O(1)=4(1-\rho)T+O(1),
\]
and the variances of $D_j$ and $L_j$ satisfy the same estimates. Since $4\rho=2\as-4$ and
$k_n=\as T+O(\log T)$, the number $s_j=k_n-d_j$ satisfies
\[
s_j=(4-\as)T+O(\log T),\qquad 2s_j-1-\E L_j=O(\log T),
\]
because $2(4-\as)=4(1-\rho)$. Thus $d_j$ and $2s_j-1$ are within $o(\sqrt T)$ of the means of $D_j$ and $L_j$, whose
variances are of order $T$. By Lemma~\ref{lem:local}, applied along an arbitrary sequence of pairs $(n,j)$ with
$n\to\infty$ and $N\le j\le2N$, there is $c>0$ such that for all large $n$ and all such $j$,
\[
\Prob(D_j=d_j)\ge\frac{c}{\sqrt T},\qquad\Prob(L_j=2s_j-1)\ge\frac{c}{\sqrt T} .
\]
Moreover $\binom{2s}{s}4^{-s}\ge1/(2\sqrt s)$ for $s\ge1$, so the last factor in~\eqref{eq:coef2} is at least
$c/\sqrt T$; also $\sum_{j=N}^{2N}1/(j+4)\ge1/3$ for large $N$, and $F_n(4)=n(n+1)(n+2)(n+3)/24\ge n^4/24$. Hence
\[
[z^{k_n}w^{k_n}]R_n\ \ge\ c\,n^4\,4^{-k_n}\,T^{-3/2}.
\]
The left side is the expected number of ordered pairs of distinct vertices which are both at depth exactly $k_n$, and
is at most $\E[Y(Y-1)]$. By~\eqref{eq:square} and the identities after it,
\[
\E\bigl[N_{n,k_n}^2\bigr]\ \ge\ n\,\E[Y(Y-1)]\ \ge\ c\,n^5\,4^{-k_n}\,T^{-3/2}
\ \ge\ \frac c4\,n^{5-\as\log4}\,T^{\frac{\log4}{2\theta}-\frac32},
\]
where the last step uses $k_n\le b_n+1$ and $4^{-b_n}=n^{-\as\log4}\,T^{(\log4)/(2\theta)}$. Since $\as<3.5912$ and
$\log4<1.3863$, we have $\as\log4<4.9785<5$. The right side therefore tends to infinity, while $\E N_{n,k_n}$ is
bounded by Lemma~\ref{lem:finite}.
\end{proof}

\begin{remark}
The proof shows which configurations are responsible. In the tree of a root $i$, the pairs kept in the proof are
two vertices at depth $k_n$ whose paths from $i$ agree up to depth about $(2\as-4)\log n\approx3.18\log n$ and then
continue separately for about $(4-\as)\log n\approx0.41\log n$ edges each. The tree of a given root reaches depth
$k_n$ with probability at most $\E Y=O(1/n)$, by Lemma~\ref{lem:finite}, whereas
$\E[Y(Y-1)]\ge c\,n^{4-\as\log4}(\log n)^{(\log4)/(2\theta)-3/2}$: a tree that reaches depth $k_n$ carries on average
at least of order $n^{5-\as\log4}$ pairs of vertices at depth $k_n$, up to a power of $\log n$.
\end{remark}

\begin{remark}
By the Cauchy--Schwarz inequality, $\Prob(M_n\ge k)\ge(\E N_{n,k})^2/\E[N_{n,k}^2]$. Proposition~\ref{prop:second}
shows that this bound tends to zero at $k=\lceil b_n\rceil$, whatever the value of $q_d$. A second moment argument at
this level would have to count a thinner family of paths, as in the truncated second moment arguments for the extremes
of branching random walks~\cite{ABR,Aidekon} and in~\cite{ABL,ABF}.
\end{remark}

\section{The top of the range, and simulations}\label{sec:top}

\begin{proof}[Proof of Proposition~\ref{prop:top}]
The coefficient of $z^{n-1}$ in~\eqref{eq:F} is $1/(n-1)!$, so Lemma~\ref{lem:first} gives
$\E N_{n,n-1}=n/(2\,(n-1)!)$. It remains to show that $N_{n,n-1}\le1$ almost surely. Suppose that the
minimum-weight path between $i$ and $j$ has $n-1$ edges; it is then a Hamilton path $\gamma$. A subpath of a
minimum-weight path is a minimum-weight path between its endpoints, since a lighter path between them would yield a
lighter walk, and hence a lighter path, between $i$ and $j$. By uniqueness, the minimum-weight path between any two
vertices is the piece of $\gamma$ between them, and it has $n-1$ edges only for the pair $\{i,j\}$. Therefore
$N_{n,n-1}$ is the indicator of $\{M_n=n-1\}$.
\end{proof}

For $n=3$ the proposition gives the law of $M_3$: it equals $2$ with probability $3/4$ and $1$ with probability
$1/4$. In general $\Prob(M_n=n-1)$ is the probability that for some vertex the tree of minimum-weight paths from it
is a path with that vertex at one end, and the proof shows that this then happens for exactly two vertices.

\subsection*{Finite checks and simulations}
The identities of Sections~\ref{sec:profile} and~\ref{sec:second} were checked by machine in small cases. The
enumeration of all recursive trees with at most $8$ vertices reproduces~\eqref{eq:F} and~\eqref{eq:coef} in exact
arithmetic, and for $K_3,\dots,K_6$ the values of $\E N_{n,k}$ in Lemma~\ref{lem:first} and of $\Prob(M_n=n-1)$
agree with estimates from $10^5$ sampled graphs within the sampling error (for $n=5$, $0.1050$ against $5/48=0.1042$).

Table~\ref{tab:sim} reports simulations of $M_n$ for $n$ between $250$ and $4000$. Each graph was sampled with
independent exponential weights and the minimum-weight paths were found by Dijkstra's algorithm from every vertex;
$R_n$ denotes the largest hopcount from vertex~$1$, which is the height of a random recursive tree. For the
comparison with the first moment let $k^*_n$ be the real number at which $\E N_{n,k}$, computed exactly
from~\eqref{eq:F} and interpolated linearly in $\log\E N_{n,k}$ between consecutive integers, equals one. By
Theorem~\ref{thm:first}, $b_n-k^*_n\to\log(1/C_*)/\theta=3.52$; for the values of $n$ in the table it lies between
$3.41$ and $3.45$.

\begin{table}[ht]
\centering
\small
\begin{tabular}{rrcccccccc}
\hline
$n$ & graphs & mean $M_n$ & s.e. & s.d. & range & mean $R_n$ & $b_n$ & $k^*_n$ & mean $M_n-k^*_n$ \\
\hline
$250$ & $8000$ & $14.03$ & $0.01$ & $1.24$ & $11$--$24$ & $10.62$ & $19.16$ & $15.75$ & $-1.72$ \\
$500$ & $4000$ & $16.31$ & $0.02$ & $1.26$ & $13$--$22$ & $12.25$ & $21.60$ & $18.18$ & $-1.86$ \\
$1000$ & $2000$ & $18.66$ & $0.03$ & $1.31$ & $15$--$25$ & $13.90$ & $24.05$ & $20.61$ & $-1.95$ \\
$2000$ & $400$ & $20.95$ & $0.06$ & $1.27$ & $18$--$25$ & $15.74$ & $26.50$ & $23.06$ & $-2.11$ \\
$4000$ & $120$ & $23.36$ & $0.10$ & $1.13$ & $21$--$28$ & $17.13$ & $28.96$ & $25.50$ & $-2.15$ \\
\hline
\end{tabular}
\caption{Simulations of $M_n$ and of $R_n$, the largest hopcount from one vertex.}\label{tab:sim}
\end{table}

\begin{table}[ht]
\centering
\small
\begin{tabular}{rccccccccc}
\hline
$n$ & $k$ & $\E N_{n,k}$ & freq. & $k$ & $\E N_{n,k}$ & freq. & $k$ & $\E N_{n,k}$ & freq. \\
\hline
$250$ & $16$ & $0.7074$ & $0.1128$ & $17$ & $0.1637$ & $0.0352$ & $18$ & $0.0353$ & $0.0083$ \\
$500$ & $19$ & $0.3148$ & $0.0505$ & $20$ & $0.0726$ & $0.0143$ & $21$ & $0.0158$ & $0.0030$ \\
$1000$ & $21$ & $0.5897$ & $0.0770$ & $22$ & $0.1429$ & $0.0290$ & $23$ & $0.0329$ & $0.0100$ \\
$2000$ & $24$ & $0.2736$ & $0.0325$ & $25$ & $0.0658$ & $0.0075$ & $26$ & $0.0151$ & $0.0000$ \\
$4000$ & $26$ & $0.5133$ & $0.0417$ & $27$ & $0.1279$ & $0.0083$ & $28$ & $0.0306$ & $0.0083$ \\
\hline
\end{tabular}
\caption{The exact value of $\E N_{n,k}$ and the observed frequency of $\{M_n\ge k\}$ at the three integers $k$
above $k^*_n$.}\label{tab:freq}
\end{table}

The mean of $M_n$ lies between $1.7$ and $2.2$ below $k^*_n$, and so between $5.1$ and $5.6$ below $b_n$. The
distance from the mean to $b_n$ grows over the range of the table: it is larger at $n=4000$ than at $n=250$ by
$0.47$, with standard error $0.10$, while $(\log\log n)/(2\theta)$ changes by $0.16$. The standard deviation of
$M_n$ is between $1.1$ and $1.3$ for every $n$. Table~\ref{tab:freq} compares the frequency of $\{M_n\ge k\}$ with
$\E N_{n,k}$, which bounds it from above. At the first integer above $k^*_n$ the frequency is $0.16$, $0.16$,
$0.13$, $0.12$ and $0.08$ of $\E N_{n,k}$ for the five values of $n$, and at the other levels in the table it is at
most $0.30$ of it. A ratio well below one is to be expected: a minimum-weight path with $k+j$ edges contains
$(j+1)(j+2)/2$ minimum-weight paths with at least $k$ edges, so that the pairs counted by $N_{n,k}$ come in groups.
The range of $n$ is too small to show whether $M_n-b_n$ is tight.

One of the $8000$ graphs with $n=250$ has $M_n=24$, while $\Prob(M_{250}\ge24)\le\E N_{250,24}=9.3\cdot10^{-7}$, so
that a sample of this size contains such a graph with probability at most $0.0075$. Its longest minimum-weight path
was recomputed by the Bellman--Ford recursion over paths with a bounded number of edges (program
\texttt{check\_rep.c}): the path has $24$ edges and weight $0.0415$, and every path with at most $23$ edges between
its endpoints is heavier. No other graph with $n=250$ has $M_n>20$.

For a single root the comparison with the first moment can be made exactly. In a random recursive tree with
$n\ge2$ vertices the subtree of the second vertex has $K$ vertices, where $K$ is uniform on $\{1,\dots,n-1\}$ because
the two sides of the tree grow as a P\'olya urn started with one ball of each colour, and, given $K$, this subtree
and the rest of the tree are independent random recursive trees with $K$ and $n-K$ vertices. Hence
$t_n(h)=\Prob(R_n>h)$ satisfies $t_n(-1)=1$, $t_1(h)=0$ for $h\ge0$ and
\[
t_n(h)=\frac1{n-1}\sum_{K=1}^{n-1}\Bigl(t_K(h-1)\bigl(1-t_{n-K}(h)\bigr)+t_{n-K}(h)\Bigr),\qquad n\ge2,\ h\ge0 .
\]
The means of $R_n$ computed from this recursion are $10.60$, $12.25$, $13.93$, $15.63$ and $17.34$ for the five
values of $n$ in Table~\ref{tab:sim}, and the simulated means differ from them by at most $1.5$ standard errors. At $k=\lceil b_n\rceil$ the ratio of $\Prob(R_n\ge k)$ to the expected number
$\sum_{d\ge k}f_{n,d}$ of vertices at depth at least $k$ is $0.62$, $0.56$, $0.51$, $0.49$ and $0.49$ for
$n=250$, $1000$, $4000$, $16000$ and $32000$.


\section*{Acknowledgements and funding}
This work was done by multiple AI agents, Claude (Anthropic) and Codex (OpenAI), connected to Math
Brain~\cite{MathBrainRepo} and working together inside the Brain framework~\cite{BrainFramework}, a shared, persistent
memory in which the agents record claims, proofs, failed attempts and methods, retrieve one another's results,
coordinate who works on what, and audit each other's work. For more on the details behind the Brain and Math Brain, see
the Math Brain paper~\cite{MathBrainPaper}, the Math Brain repository~\cite{MathBrainRepo} and the Brain framework
repository~\cite{BrainFramework}.

The research presented in this paper was supported by the European Research Council (ERC) under the European Union's
Horizon Europe research and innovation programme (grant agreement No.~101041711), by the Simons Foundation as part of
the Collaboration on the Mathematical and Scientific Foundations of Deep Learning, by Heights Labs, by Convex Nexus
Capital, by the Israel Science Foundation (grant number 2258/19), and by the Israel Science Foundation (ISF Grant
4101/25). Computational resources were provided in part by Google Cloud through its education credits program (credit
530550103).

\section*{Declaration of competing interest}
The author holds equity in Heights Labs and in Convex Nexus Capital and has a management role at Convex Nexus Capital.
These entities are included in the funding declaration above.

\section*{Declaration of generative AI and AI-assisted technologies}
The paper comes out of Math Brain, an ongoing project on the use of AI systems for research mathematics, and its
results were developed with the Brain~\cite{BrainFramework} and Math Brain~\cite{MathBrainRepo}. In October 2026 two
AI agents connected to the Brain, working independently from the project's list of open problems in probability,
obtained Theorems~\ref{thm:upper} and~\ref{thm:first} and Propositions~\ref{prop:second} and~\ref{prop:top} with
proofs, and a third agent checked the arguments. In November 2026 a Claude Fable agent (Anthropic) checked the proofs
again, wrote and ran the programs and prepared the manuscript. These checks are not peer review. The author directed
the work and is responsible for the content of the paper and for decisions concerning its dissemination.

\section*{Code and data availability}
Programs and outputs:\\
\url{https://github.com/yspennstate/longest-geodesic-complete-graph}.

The paper is part of the Math Brain project, whose public repository is
\url{https://github.com/yspennstate/math-brain-public}~\cite{MathBrainRepo}.


\begin{thebibliography}{99}
\bibitem{ABL} L.~Addario-Berry, N.~Broutin and G.~Lugosi, The longest minimum-weight path in a complete graph,
\emph{Combin. Probab. Comput.} \textbf{19} (2010), 1--19. \href{https://doi.org/10.1017/S0963548309990204}{doi:10.1017/S0963548309990204}

\bibitem{ABF} L.~Addario-Berry and K.~Ford, Poisson--Dirichlet branching random walks,
\emph{Ann. Appl. Probab.} \textbf{23} (2013), 283--307. \href{https://doi.org/10.1214/12-AAP840}{doi:10.1214/12-AAP840}

\bibitem{ABR} L.~Addario-Berry and B.~Reed, Minima in branching random walks,
\emph{Ann. Probab.} \textbf{37} (2009), 1044--1079. \href{https://doi.org/10.1214/08-AOP428}{doi:10.1214/08-AOP428}

\bibitem{Aidekon} E.~A\"{\i}d\'ekon, Convergence in law of the minimum of a branching random walk,
\emph{Ann. Probab.} \textbf{41} (2013), 1362--1426. \href{https://doi.org/10.1214/12-AOP750}{doi:10.1214/12-AOP750}

\bibitem{BvdH} S.~Bhamidi and R.~van der Hofstad, Diameter of the stochastic mean-field model of distance,
\emph{Combin. Probab. Comput.} \textbf{26} (2017), 797--825. \href{https://doi.org/10.1017/S0963548317000232}{doi:10.1017/S0963548317000232}

\bibitem{ChenMa} X.~Chen and H.~Ma, Atypical decay rates for atypical heights in random recursive trees,
preprint, 2026. \href{https://arxiv.org/abs/2604.20139}{arXiv:2604.20139}

\bibitem{Devroye} L.~Devroye, Branching processes in the analysis of the heights of trees,
\emph{Acta Inform.} \textbf{24} (1987), 277--298. \href{https://doi.org/10.1007/BF00265991}{doi:10.1007/BF00265991}

\bibitem{Devroye88} L.~Devroye, Applications of the theory of records in the study of random trees,
\emph{Acta Inform.} \textbf{26} (1988), 123--130. \href{https://doi.org/10.1007/BF02915448}{doi:10.1007/BF02915448}

\bibitem{Drmota} M.~Drmota, \emph{Random Trees: An Interplay between Combinatorics and Probability}, Springer, Vienna, 2009.
\href{https://doi.org/10.1007/978-3-211-75357-6}{doi:10.1007/978-3-211-75357-6}

\bibitem{EGHN1} M.~Eckhoff, J.~Goodman, R.~van der Hofstad and F.~R.~Nardi, Long paths in first passage percolation on
the complete graph I. Local PWIT dynamics, \emph{Electron. J. Probab.} \textbf{25} (2020), 1--45.
\href{https://doi.org/10.1214/20-EJP484}{doi:10.1214/20-EJP484}

\bibitem{EGHN2} M.~Eckhoff, J.~Goodman, R.~van der Hofstad and F.~R.~Nardi, Long paths in first passage percolation on
the complete graph II. Global branching dynamics, \emph{J. Stat. Phys.} \textbf{181} (2020), 364--447.
\href{https://doi.org/10.1007/s10955-020-02585-1}{doi:10.1007/s10955-020-02585-1}

\bibitem{FHN} M.~Fuchs, H.-K.~Hwang and R.~Neininger, Profiles of random trees: limit theorems for random recursive
trees and binary search trees, \emph{Algorithmica} \textbf{46} (2006), 367--407.
\href{https://doi.org/10.1007/s00453-006-0109-5}{doi:10.1007/s00453-006-0109-5}

\bibitem{vdH} R.~van der Hofstad, \emph{Stochastic processes on random graphs}, lecture notes for the 47th Summer
School in Probability, Saint-Flour, 2017.
\url{https://www.win.tue.nl/~rhofstad/SaintFlour_SPoRG_fin_Saint_Flour.pdf}

\bibitem{vdHHVM} R.~van der Hofstad, G.~Hooghiemstra and P.~Van Mieghem, First-passage percolation on the random
graph, \emph{Probab. Engrg. Inform. Sci.} \textbf{15} (2001), 225--237.
\href{https://doi.org/10.1017/S026996480115206X}{doi:10.1017/S026996480115206X}

\bibitem{HVM} G.~Hooghiemstra and P.~Van Mieghem, The weight and hopcount of the shortest path in the complete graph
with exponential weights, \emph{Combin. Probab. Comput.} \textbf{17} (2008), 537--548.
\href{https://doi.org/10.1017/S0963548308009176}{doi:10.1017/S0963548308009176}

\bibitem{Hwang} H.-K.~Hwang, Asymptotic expansions for the Stirling numbers of the first kind,
\emph{J. Combin. Theory Ser. A} \textbf{71} (1995), 343--351.
\href{https://doi.org/10.1016/0097-3165(95)90010-1}{doi:10.1016/0097-3165(95)90010-1}

\bibitem{Janson} S.~Janson, One, two and three times $\log n/n$ for paths in a complete graph with random weights,
\emph{Combin. Probab. Comput.} \textbf{8} (1999), 347--361.
\href{https://doi.org/10.1017/S0963548399003892}{doi:10.1017/S0963548399003892}

\bibitem{MW} L.~Moser and M.~Wyman, Asymptotic development of the Stirling numbers of the first kind,
\emph{J. London Math. Soc.} \textbf{33} (1958), 133--146.
\href{https://doi.org/10.1112/jlms/s1-33.2.133}{doi:10.1112/jlms/s1-33.2.133}

\bibitem{PS1} M.~Pain and D.~S\'enizergues, Correction terms for the height of weighted recursive trees,
\emph{Ann. Appl. Probab.} \textbf{32} (2022), 3027--3059. \href{https://doi.org/10.1214/21-AAP1756}{doi:10.1214/21-AAP1756}

\bibitem{PS2} M.~Pain and D.~S\'enizergues, Height of weighted recursive trees with sub-polynomially growing total
weight, \emph{Ann. Inst. Henri Poincar\'e Probab. Stat.} \textbf{60} (2024), 1949--1984.
\href{https://doi.org/10.1214/23-AIHP1379}{doi:10.1214/23-AIHP1379}

\bibitem{Pittel} B.~Pittel, Note on the heights of random recursive trees and random $m$-ary search trees,
\emph{Random Structures Algorithms} \textbf{5} (1994), 337--347.
\href{https://doi.org/10.1002/rsa.3240050207}{doi:10.1002/rsa.3240050207}

\bibitem{SM} R.~T.~Smythe and H.~M.~Mahmoud, A survey of recursive trees,
\emph{Theory Probab. Math. Statist.} \textbf{51} (1995), 1--27.

\bibitem{MathBrainRepo} Y.~Shmalo, \emph{Math Brain: public snapshot, code and tests}, 2026,
\url{https://github.com/yspennstate/math-brain-public}.
\bibitem{MathBrainPaper} Y.~Shmalo, \emph{Math Brain: shared, auditable memory for multi-agent mathematical research},
October 2026, \url{https://github.com/yspennstate/math-brain-public/tree/main/papers/math-brain-2026-10-24}.
\bibitem{BrainFramework} Y.~Shmalo, \emph{Brain Framework}, public software repository, 2026,
\url{https://github.com/yspennstate/brain-framework-public}.

\end{thebibliography}
\end{document}
